\documentclass{article}
\usepackage{amsmath, amssymb, amsthm}
\usepackage{hyperref}
\usepackage{geometry}
\geometry{margin=1in}

\title{Erd\H{o}s Problem \#104}
\date{Last updated: 2025-08-31}
\author{Source: erdosproblems.com}

\begin{document}
\maketitle

\noindent\textbf{Status:} OPEN \quad \textbf{Prize: \$100}

\noindent\textbf{Tags:} geometry

\medskip

\noindent\textbf{Problem Statement:}

Given $n$ points in $\mathbb{R}^2$ the number of distinct unit circles containing at least three points is $o(n^2)$.




In \cite{Er81d} Erd\H{o}s proved that $\gg n$ many circles is possible, and that there cannot be more than $O(n^2)$ many circles. The argument is very simple: every pair of points determines at most $2$ unit circles, and the claimed bound follows from double counting. Erd\H{o}s claims in a number of places this produces the upper bound $n(n-1)$, but Harborth and Mengerson \cite{HaMe86} note that in fact this delivers an upper bound of $\frac{n(n-1)}{3}$. Elekes \cite{El84} has a simple construction of a set with $\gg n^{3/2}$ such circles. This may be the correct order of magnitude. In \cite{Er75h} and \cite{Er92e} Erd\H{o}s also asks how many such unit circles there must be if the points are in general position. In \cite{Er92e} Erd\H{o}s offered £100 for a proof or disproof that the answer is $O(n^{3/2})$. The maximal number of unit circles achieved by $n$ points is A003829 in the OEIS. See also [506] and [831] .




References

[El84] Elekes, G., {$n$} points in the plane can determine $n^{3/2}$ unit circles . Combinatorica (1984), 131.

[Er75h] Erd\H{o}s, P., Some problems on elementary geometry . Austral. Math. Soc. Gaz. (1975), 2-3.

[Er81d] Erd\H{o}s, P., Some applications of graph theory and combinatorial methods to number theory and geometry . Algebraic methods in graph theory, Vol. I, II (Szeged, 1978) (1981), 137-148.

[Er92e] Erd\H{o}s, P\'{a}l, Some Unsolved problems in Geometry, Number Theory and Combinatorics . Eureka (1992), 44-48.

[HaMe86] Harborth, Heiko and Mengersen, Ingrid, Point sets with many unit circles . Discrete Math. (1986), 193--197.

\noindent\textbf{Related OEIS sequences:} A003829


\bigskip
\noindent\small{Source: \url{https://www.erdosproblems.com/104}}
\end{document}
